% 颜色是全书和 TikZ 图形共享的设计变量。
\definecolor{BookInk}{HTML}{233342}
\definecolor{BookTeal}{HTML}{267D82}
\definecolor{BookBlue}{HTML}{42658A}
\definecolor{BookAmber}{HTML}{A66A25}
\definecolor{BookMuted}{HTML}{64717B}
\definecolor{BookPaper}{HTML}{F3F5F5}
\definecolor{BookRule}{HTML}{DCE3E5}
% 科研示意图：白底、深灰细线与柔和填充；与封面和正文调色板分开维护。
\definecolor{FigureInk}{HTML}{222222}
\definecolor{FigureStroke}{HTML}{444444}
\definecolor{FigureMuted}{HTML}{686868}
\definecolor{FigureGrid}{HTML}{C4C9CC}
\definecolor{FigurePaper}{HTML}{F4F5F6}
% 离散强调色按种类数选择：单色黄、双色蓝红、三色蓝红黄。
\definecolor{FigureBlue}{HTML}{7998AD}
\definecolor{FigureRed}{HTML}{D57B70}
\definecolor{FigureYellow}{HTML}{D6B35D}
\definecolor{FigureBlueFill}{HTML}{EDF2F5}
\definecolor{FigureRedFill}{HTML}{F8E8E5}
\definecolor{FigureYellowFill}{HTML}{FFF1CF}
\definecolor{FigureCream}{HTML}{FBEFDA}
% 其他卷保留原有默认色；第二卷在局部分组中覆盖。
\colorlet{FigureInput}{BookBlue}
\colorlet{FigureModel}{BookTeal}
\colorlet{FigureOutput}{BookInk}
\colorlet{FigureLoss}{BookAmber}
% 神经网络图：白色画布与浅色节点保留拓扑，明亮强调色和清晰线条标出计算职责。
\definecolor{NNInput}{HTML}{2687EA}
\definecolor{NNHidden}{HTML}{9464DA}
\definecolor{NNOutput}{HTML}{12ADA7}
\definecolor{NNGradient}{HTML}{A95B13}
\definecolor{NNLine}{HTML}{8A9CB6}
\colorlet{NNInk}{BookInk}
\colorlet{NNCanvas}{white}
\colorlet{NNPanel}{NNInput!4!white}
\definecolor{NNWire}{HTML}{657B99}
\colorlet{FigureInputFill}{FigureInput!22!white}
\colorlet{FigureModelFill}{FigureModel!18!white}
\colorlet{FigureOutputFill}{FigureOutput!20!white}
\colorlet{FigureLossFill}{FigureLoss!18!white}
\colorlet{FigurePanel}{FigureModel!9!white}
\colorlet{FoundationInput}{FigureInput}
\colorlet{FoundationModel}{FigureModel}
\colorlet{FoundationOutput}{FigureOutput}
\colorlet{FoundationLoss}{FigureLoss}

\definecolor{BookTheorem}{HTML}{2F65B0}
\colorlet{BookLemma}{BookTheorem!60!white}
\definecolor{BookExample}{HTML}{D3AD38}
\definecolor{BookExampleBackground}{HTML}{FFF8D9}

% Letter 纸张：正文 112 mm，边注 48 mm，中缝 9 mm。
\geometry{left=23mm,top=24mm,textwidth=112mm,textheight=226mm,
  marginparsep=9mm,marginparwidth=48mm,headheight=23pt,headsep=9mm}
\renewcommand{\contentsname}{目录}
\renewcommand{\figurename}{图}
\renewcommand{\tablename}{表}
\renewcommand{\thepart}{\arabic{part}}
\newcommand{\BookChapterName}{第\arabic{chapter}章}
\newcommand{\BookPartName}{第\chinese{part}部分}
\renewcommand{\allcaps}[1]{\MakeUppercase{#1}}
\renewcommand{\smallcaps}[1]{#1}
\setcounter{secnumdepth}{3}
\setcounter{tocdepth}{1}
\setlength{\parindent}{2em}
\setlength{\parskip}{0.35\baselineskip plus 1pt}
\linespread{1.16}
\emergencystretch=1em
\raggedbottom
% 跨页段落至少保留两行，避免新页只剩一句的尾字。
\clubpenalty=10000
\widowpenalty=10000
\displaywidowpenalty=10000
\makeatletter
\@clubpenalty=10000
\makeatother

% 跨页表沿用全版心和正文表格字号；各页表头由内容源声明。
\newenvironment{BookLongTable}[1]{%
  \begin{fullwidth}\small
  \setlength{\hsize}{\linewidth}%
  \setlength{\parindent}{0pt}%
  \setlength{\RaggedRightParindent}{0pt}%
  % 完整读取当前问题地图，避免默认分块切断跨行组的续页表头。
  \setcounter{LTchunksize}{1000}%
  \renewcommand{\multirowsetup}{\RaggedRight\setlength{\parindent}{0pt}}%
  \booktabletopcells
  \renewcommand{\arraystretch}{1.15}%
  \setlength{\LTleft}{0pt}\setlength{\LTright}{\fill}%
  \setlength{\LTcapwidth}{\linewidth}%
  \longtable{#1}%
}{\endlongtable\end{fullwidth}}
\makeatletter
% 顶部对齐的合并单元格按内容的自然高度设置。
% multirow 的默认高度按单行估计，不适合本表实际占多行的段落单元格。
\newcommand{\booktabletopcells}{%
  \let\booktable@multirow@vbox\multirow@vbox
  \long\def\multirow@vbox##1##2##3{%
    \if ##1t\relax
      \setbox\z@\vtop{##2\multirowsetup ##3}%
    \else
      \booktable@multirow@vbox{##1}{##2}{##3}%
    \fi}}
% longtable 使用自己的图表注入口，保持与普通表注相同的字体、颜色和左对齐。
\long\def\LT@makecaption#1#2#3{%
  \LT@mcol\LT@cols c{\hbox to\z@{\hss\parbox[t]\LTcapwidth{%
    \booknotefont\RaggedRight
    #1{{\color{BookTeal}#2}\quad}#3\par\vskip\baselineskip}%
  \hss}}}
\makeatother

\newcommand{\booknotefont}{\kaishu\fontsize{9}{12.5}\selectfont}
\setsidenotefont{\booknotefont}
\setmarginnotefont{\booknotefont}
\setcaptionfont{\booknotefont}
\AddToHook{env/quote/begin}{\kaishu}
\AddToHook{env/quotation/begin}{\kaishu}
\makeatletter
\renewcommand{\@tufte@sidenote@justification}{\RaggedRight}
\renewcommand{\@tufte@marginnote@justification}{\RaggedRight}
\makeatother

\titleformat{\chapter}[display]
  {\heiti\color{BookInk}}
  {\colorbox{BookTeal}{\parbox[c][10mm][c]{28mm}{\centering
    \color{white}\fontsize{16}{20}\selectfont\BookChapterName}}}
  {5mm}{\fontsize{27}{34}\selectfont\bfseries}
\titleformat{name=\chapter,numberless}[display]
  {\heiti\color{BookInk}}
  {\color{BookTeal}\rule{28mm}{3mm}}
  {5mm}{\fontsize{27}{34}\selectfont\bfseries}
\titlespacing*{\chapter}{0pt}{6mm}{12mm}
\titleformat{\section}{\heiti\Large\bfseries\color{BookInk}}{\thesection}{0.7em}{}
\titleformat{\subsection}{\heiti\large\bfseries\color{BookInk}}{\thesubsection}{0.7em}{}
% Tufte 禁用此层级；在项目样式层恢复，保留模型章节的原有结构。
\titleclass{\subsubsection}{straight}
\titleformat{\subsubsection}{\heiti\normalsize\bfseries\color{BookInk}}{\thesubsubsection}{0.7em}{}
\titleformat{\paragraph}[runin]{\heiti\bfseries\color{BookInk}}{}{0pt}{}
\titlespacing*{\section}{0pt}{2.5ex plus 1ex}{1ex}
\titlespacing*{\subsection}{0pt}{2ex plus 1ex}{0.8ex}
\titlespacing*{\subsubsection}{0pt}{1.8ex plus 0.8ex}{0.6ex}

\makeatletter
\newcommand{\bookpagebadge}{%
  \begingroup\setlength{\fboxsep}{4pt}%
  \colorbox{BookTeal}{\makebox[6mm]{\sffamily\small\color{white}\thepage}}%
  \endgroup}
\newcommand{\bookchapterrunningtitle}{%
  \heiti\footnotesize\color{BookMuted}\nouppercase{\leftmark}}
\newcommand{\booksectionrunningtitle}{%
  \heiti\footnotesize\color{BookMuted}\nouppercase{\rightmark}}
\newcommand{\bookoddpagerunningtitle}{%
  \if@mainmatter\booksectionrunningtitle\else\bookchapterrunningtitle\fi}
\fancypagestyle{book}{%
  \fancyhf{}
  % 偶数页显示章标题，奇数正文页显示节标题；页码始终位于外侧。
  \fancyhead[L]{\ifodd\value{page}\bookoddpagerunningtitle\else\bookpagebadge\fi}
  \fancyhead[R]{\ifodd\value{page}\bookpagebadge\else\bookchapterrunningtitle\fi}
  \renewcommand{\headrulewidth}{0pt}}
\fancypagestyle{plain}{\fancyhf{}}
\renewcommand{\chaptermark}[1]{%
  \markboth{\if@mainmatter\BookChapterName\quad\fi #1}{}}
\renewcommand{\sectionmark}[1]{%
  \markright{\if@mainmatter\thesection\quad\fi #1}}
\renewcommand{\frontmatter}{\clearpage\@mainmatterfalse\pagenumbering{roman}\pagestyle{book}}
\renewcommand{\mainmatter}{%
  \clearpage
  % 封面与罗马页码各自重新编号；用实际输出页数保证正文第 1 页在右页。
  \ifodd\numexpr\ReadonlyShipoutCounter+1\relax\else
    \null\thispagestyle{empty}\newpage
  \fi
  \@mainmattertrue\pagenumbering{arabic}\pagestyle{book}}

% 恢复常规浮动体，取消 Tufte 将图表注自动移到边栏的行为。
\renewenvironment{figure}[1][htbp]{\@float{figure}[#1]}{\end@float}
\renewenvironment{table}[1][htbp]{\@float{table}[#1]}{\end@float}
\renewenvironment{figure*}[1][htbp]{\@float{figure}[#1]\begin{fullwidth}}{\end{fullwidth}\end@float}
\renewenvironment{table*}[1][htbp]{\@float{table}[#1]\begin{fullwidth}}{\end{fullwidth}\end@float}
\long\def\@caption#1[#2]#3{%
  \hyper@makecurrent{#1}\NR@gettitle{#2}%
  \par\addcontentsline{\csname ext@#1\endcsname}{#1}%
    {\protect\numberline{\csname the#1\endcsname}{\ignorespaces #2}}%
  \begingroup\@parboxrestore
    \if@minipage\@setminipage\fi
    \vskip7pt\booknotefont\RaggedRight
    \noindent\Hy@raisedlink{\hyper@anchorstart{\@currentHref}\hyper@anchorend}%
    {\color{BookTeal}\csname fnum@#1\endcsname}\quad
    \ignorespaces#3\par
  \endgroup}
\makeatother

% 双语词条需要更宽的栏宽；在项目层覆盖 Tufte 默认的三栏索引。
\makeatletter
\renewenvironment{theindex}{%
  \chapter{\indexname}%
  \begin{fullwidth}%
  \small
  \parskip0pt
  \parindent0pt
  \let\item\@idxitem
  \begin{multicols}{2}%
}{%
  \end{multicols}%
  \end{fullwidth}%
}
\makeatother

% 草稿章节可以先保留图位；资源落地后自动换成正式图片。
% 缺图时在 PDF 中明确显示文件名，避免静默遗漏。
\NewDocumentCommand{\bookdraftgraphic}{O{width=\linewidth} m}{%
  \IfFileExists{#2}
    {\includegraphics[#1]{#2}}
    {\PackageWarning{machine-learning}{Draft figure `#2' is not available}%
     \begingroup
       \setlength{\fboxsep}{0pt}%
       \fcolorbox{BookRule}{BookPaper}{%
         \parbox[c][42mm][c]{\dimexpr\linewidth-2\fboxrule\relax}{%
           \centering\heiti\color{BookMuted}\small
           配图待补充\par\medskip
           \ttfamily\footnotesize\detokenize{#2}}}%
     \endgroup}}

% 前置页均是独立版面；只在这些版面命令内使用显式定位。
\newcommand{\makebookcover}{%
  \thispagestyle{empty}\null
  \begin{tikzpicture}[remember picture,overlay]
    \fill[BookInk] (current page.north west) rectangle (current page.south east);
    \fill[BookTeal] ([xshift=23mm,yshift=-27mm]current page.north west) rectangle ++(17mm,-7mm);
    \node[anchor=north west,text=white,font=\sffamily\bfseries\fontsize{11}{15}\selectfont]
      at ([xshift=23mm,yshift=-42mm]current page.north west) {\BookEnglishTitle};
    \node[anchor=north west,inner sep=0pt,text=white,
      font=\heiti\bfseries\fontsize{34}{42}\selectfont] (book-cover-title)
      at ([xshift=23mm,yshift=-68mm]current page.north west) {\BookTitle};
    \node[anchor=north west,inner sep=0pt,text=white!75!BookInk,
      font=\heiti\mdseries\fontsize{18}{24}\selectfont]
      at ([yshift=-6mm]book-cover-title.south west) {\BookCoverSecondLine};
    \fill[BookTeal] ([xshift=138mm,yshift=38mm]current page.south west) rectangle ++(54mm,54mm);
    \fill[white!85!BookTeal] ([xshift=168mm,yshift=38mm]current page.south west) rectangle ++(24mm,24mm);
    \node[anchor=south west,text=white,font=\heiti\mdseries\fontsize{13}{18}\selectfont]
      at ([xshift=23mm,yshift=38mm]current page.south west) {\BookAuthor};
    \node[anchor=south west,text=white!70!BookInk,font=\heiti\small]
      at ([xshift=23mm,yshift=30mm]current page.south west) {\BookEditionLabel};
  \end{tikzpicture}\clearpage}

\newcommand{\makebooktitle}{%
  \thispagestyle{empty}
  \begin{fullwidth}
    \setlength{\parindent}{0pt}
    \vspace*{22mm}
    {\color{BookTeal}\rule{17mm}{3mm}\par}
    \vspace{10mm}
    {\sffamily\bfseries\color{BookMuted}\fontsize{11}{15}\selectfont
      \BookEnglishTitle\par}
    \vspace{5mm}
    {\heiti\bfseries\color{BookInk}\fontsize{34}{42}\selectfont\BookTitle\par}
    \vspace{6mm}
    {\heiti\mdseries\color{BookMuted}\fontsize{18}{24}\selectfont\BookCoverSecondLine\par}
    \vspace{15mm}
    {\heiti\large\BookAuthor\par}
    \vspace{3mm}
    {\heiti\small\color{BookMuted}\BookEditionLabel\par}
    \vfill
  \end{fullwidth}\clearpage}

\newcommand{\makebookcopyright}{%
  \begingroup\setlength{\parindent}{0pt}
  \thispagestyle{empty}\null
  \begin{tikzpicture}[remember picture,overlay]
    \node[anchor=south east,inner sep=0pt,text=BookMuted!10,
      font=\fontsize{104}{104}\selectfont]
      at ([xshift=-23mm,yshift=33mm]current page.south east) {\ccLogo};
  \end{tikzpicture}
  \vfill
  {\heiti\bfseries\BookTitle\par}
  {\sffamily\small\BookEnglishTitle\par}
  \bigskip
  {\kaishu\small
    作者：\BookAuthor\par
    版本：\BookEdition\par
    年份：\BookYear\par
    Copyright \textcopyright\ \BookYear\ \BookAuthor\par}
  \medskip
  {\color{BookMuted}\fontsize{18}{22}\selectfont\ccbyncnd\par}
  \smallskip
  {\rmfamily\small\BookCopyright\par}
  \clearpage\endgroup}

% 目录中的卷使用整行青色底，部分保留原来的青色编号块。
\newlength{\booktoctitleindent}
\setlength{\booktoctitleindent}{29mm}
\newcommand{\booktocvolumeband}[1]{%
  \rlap{\color{BookTeal}\rule[-3mm]{\linewidth}{11mm}}%
  \hspace*{4mm}%
  \makebox[\dimexpr\booktoctitleindent-4mm\relax][l]{\color{white}第\zhnumber{#1}卷}}
\newcommand{\booktocpartbadge}[1]{%
  \colorbox{BookTeal}{\heiti\small\color{white}第\zhnumber{#1}部分}}
\newcommand{\booktocbackmatterbadge}[1]{%
  \colorbox{BookTeal}{\heiti\small\color{white}#1}}
% 各级条目共用页码栏，右边界距目录色块右缘 4 mm。
\newcommand{\booktocpage}{%
  \makebox[10mm][r]{\thecontentspage}\hspace*{4mm}}

% 卷首页：整页墨蓝，卷号、中英文标题组成统一左对齐的文字组。
\newcommand{\bookvolumetitlepage}[2]{%
  \begin{tikzpicture}[remember picture,overlay]
    \fill[BookInk] (current page.north west) rectangle (current page.south east);
    \node[anchor=north west,inner sep=0pt,text=white!78!BookInk,
      font=\heiti\bfseries\fontsize{18}{24}\selectfont]
      at ([xshift=23mm,yshift=-82mm]current page.north west) {\BookVolumeName};
    \node[anchor=north west,inner sep=0pt,text=white,text width=169mm,
      align=left,font=\heiti\bfseries\fontsize{38}{50}\selectfont]
      (book-volume-title)
      at ([xshift=23mm,yshift=-110mm]current page.north west) {#1};
    \node[anchor=north west,inner sep=0pt,text=white!78!BookInk,
      text width=169mm,align=left,font=\sffamily\bfseries\fontsize{13}{19}\selectfont]
      at ([yshift=-10mm]book-volume-title.south west) {#2};
  \end{tikzpicture}}

% 分部首页仅显示部分编号和中英文标题，不显示所属卷的页眉。
% 卷、部分与目录共享全书计数；新卷不重置部分号和章号。
% 目录卷标题预留与下属条目同页的空间，避免空部分组跨页散开。
\usepackage{needspace}
\newcounter{volume}
\renewcommand{\thevolume}{\arabic{volume}}
\newcommand{\BookVolumeName}{第\chinese{volume}卷}
\newcommand{\BookVolumeTitle}{}
\newcommand{\BookPartEnglishTitle}{}
\ifBookVolumeEdition
  \newcommand{\BookTopBookmarkLevel}{-1}
\else
  \newcommand{\BookTopBookmarkLevel}{-2}
\fi
\newcommand{\BookTopBookmark}[2]{%
  \pdfbookmark[\BookTopBookmarkLevel]{#1}{#2}}

\NewDocumentCommand{\BookVolumeRef}{m m}{%
  \ifBookVolumeEdition #2\else 第~\ref{#1}卷\fi}
\NewDocumentCommand{\BookPartRef}{m m}{%
  \ifBookVolumeEdition #2\else 第~\ref{#1}部分\fi}
\NewDocumentCommand{\BookChapterRef}{m m}{%
  \ifBookVolumeEdition #2\else 第~\ref{#1}章\fi}
\NewDocumentCommand{\BookSectionRef}{m m}{%
  \ifBookVolumeEdition #2\else 第~\ref{#1}节\fi}
\NewDocumentCommand{\BookTheoremRef}{m m}{%
  \ifBookVolumeEdition #2\else 定理~\ref{#1}\fi}
\NewDocumentCommand{\BookEquationRef}{m m}{%
  \ifBookVolumeEdition #2\else 式~\eqref{#1}\fi}

\makeatletter
% 注册高于 part 的目录及 PDF 书签层级，不改动第三方模板。
\def\ttll@volume{-2}
\def\toclevel@volume{-2}
\def\ttll@bookbackmatter{-1}
\def\toclevel@bookbackmatter{-1}
\providecommand*{\theHvolume}{\arabic{volume}}
\def\Hy@numberline#1{%
  \ifnum\Hy@level=-2
    第\zhnumber{#1}卷\space
  \else\ifnum\Hy@level=-1
    第\zhnumber{#1}部分\space
  \else\ifnum\Hy@level=0
    第 #1 章\space
  \else
    #1\space
  \fi\fi\fi}
\newcommand{\BookBackmatterTocEntry}[1]{%
  \begingroup
    \long\def\Hy@writebookmark##1##2##3##4##5{}%
    \addcontentsline{toc}{bookbackmatter}{\protect\booktocbackmatterbadge{#1}}%
  \endgroup}
\newcommand{\BookBackmatterChapter}[2]{%
  \chapter*{#1}%
  \BookTopBookmark{#1}{#2}%
  \markboth{#1}{}%
  \BookBackmatterTocEntry{#1}}
\newcommand{\BookBackmatterIndexChapter}[1]{%
  \BookOriginalChapter*{#1}%
  \BookTopBookmark{#1}{book-index}%
  \markboth{#1}{}%
  \BookBackmatterTocEntry{#1}}
\renewcommand{\tableofcontents}{%
  \chapter*{\contentsname}\markboth{\contentsname}{}%
  \BookTopBookmark{\contentsname}{book-contents}%
  \begin{fullwidth}
    \@starttoc{toc}%
  \end{fullwidth}}
\makeatother

\contentsmargin{0pt}
\titlecontents{volume}[0pt]
  {\Needspace{60mm}\addvspace{17pt}\contentsmargin{0pt}%
    \heiti\large\mdseries\color{white}\hypersetup{linkcolor=white}}
  {\booktocvolumeband{\thecontentslabel}}{}%
  {\hfill\booktocpage}
  [\vspace{5pt}]
\titlecontents{part}[\booktoctitleindent]
  {\addvspace{9pt}\heiti\bfseries\color{BookTeal}}
  {\contentslabel[\booktocpartbadge{\thecontentslabel}]{25mm}}{}%
  {\hspace{0.6em}\titlerule*[0.6pc]{\textcolor{BookRule}{.}}\booktocpage}
\titlecontents{bookbackmatter}[\booktoctitleindent]
  {\addvspace{9pt}\heiti\bfseries\color{BookTeal}}
  {}{\hspace*{-25mm}}%
  {\hspace{0.6em}\titlerule*[0.6pc]{\textcolor{BookRule}{.}}\booktocpage}
\titlecontents{chapter}[\booktoctitleindent]
  {\addvspace{5pt}\heiti\color{BookInk}}
  {\contentslabel[第\thecontentslabel 章]{22mm}}{}%
  {\hfill\booktocpage}
\titlecontents{section}[\booktoctitleindent]
  {\small\color{BookMuted}}
  {\contentslabel{22mm}}{}{\hfill\booktocpage}

% 卷首页只接收中英文标题；具体视觉样式留在 layout.tex 中。
\makeatletter
\newcommand{\bookvolume}[2]{%
  \ifBookVolumeEdition\else
    \clearpage\thispagestyle{empty}%
  \fi
  \refstepcounter{volume}%
  \renewcommand{\BookVolumeTitle}{#1}%
  \NR@gettitle{#1}%
  \markboth{\BookVolumeName\quad #1}{}%
  \ifBookVolumeEdition\else
    \addcontentsline{toc}{volume}{\protect\numberline{\thevolume}#1}%
    \null
    \bookvolumetitlepage{#1}{#2}\par
  \fi
  \ignorespaces}

\newcommand{\bookpart}[2]{%
  \renewcommand{\BookPartEnglishTitle}{#2}\part{#1}}
\RenewDocumentCommand{\part}{s o m}{%
  % 不依赖 openright 或单双面选项：部首页始终使用奇数书页码。
  % 在建立标签、目录与书签锚点之前补空白页，避免链接落到空白页。
  \clearpage
  \ifodd\value{page}\else
    \null\thispagestyle{empty}\newpage
  \fi
  \thispagestyle{empty}
  \IfBooleanTF{#1}{\phantomsection}{\refstepcounter{part}}
  \NR@gettitle{#3}%
  \IfBooleanF{#1}{\IfNoValueTF{#2}
    {\addcontentsline{toc}{part}{\protect\numberline{\thepart}#3}}
    {\addcontentsline{toc}{part}{\protect\numberline{\thepart}#2}}}
  \markboth{#3}{}\null
  \begin{tikzpicture}[remember picture,overlay]
    \fill[BookPaper] (current page.north west) rectangle (current page.south east);
    \fill[BookTeal] ([xshift=23mm,yshift=-54mm]current page.north west) rectangle ++(57mm,-22mm);
    \node[anchor=center,text=white,font=\heiti\fontsize{21}{27}\selectfont]
      at ([xshift=51.5mm,yshift=-65mm]current page.north west)
      {\IfBooleanF{#1}{\BookPartName}};
    \node[anchor=north west,inner sep=0pt,text=BookInk,text width=169mm,
      font=\heiti\bfseries\fontsize{29}{38}\selectfont]
      at ([xshift=23mm,yshift=-106mm]current page.north west) {#3};
    \node[anchor=north west,inner sep=0pt,text=BookMuted,
      font=\sffamily\fontsize{11}{16}\selectfont]
      at ([xshift=23mm,yshift=-126mm]current page.north west) {\BookPartEnglishTitle};
  % 到下一个结构命令才输出，使随后的 label 指向当前入口页。
  \end{tikzpicture}\par\ignorespaces}
\makeatother


% 全书图形的基础语法；局部几何与数据编码可覆盖对应参数。
\tikzset{
  figure picture/.style={font=\figurefont\footnotesize,text=FigureInk,
    line width=.6pt},
  figure node/.style={rectangle,rounded corners=1.5mm,
    draw=FigureStroke,fill=FigureModelFill,line width=.75pt,
    minimum height=10mm,inner xsep=3pt,inner ysep=4pt,
    font=\figurefont\footnotesize,align=center,text=FigureInk},
  figure flow/.style={-{Stealth[length=1.8mm,width=1.2mm]},
    draw=FigureStroke,line width=.85pt},
  figure link/.style={draw=FigureMuted,line width=.6pt},
  figure panel/.style={draw=FigureStroke,fill=white,line width=.75pt,
    rounded corners=2.5mm},
  figure subgroup/.style={draw=FigureMuted,line width=.6pt,
    rounded corners=2mm,dash pattern=on 2pt off 2pt},
  figure heading/.style={font=\figurefont\bfseries\footnotesize,text=FigureInk},
  figure note/.style={font=\figurefont\footnotesize,text=FigureMuted,align=center},
  figure edge label/.style={font=\figurefont\footnotesize,text=FigureInk,
    fill=white,inner sep=1.5pt,align=center},
  book node/.style={figure node,minimum height=10mm,inner sep=6pt,
    font=\figurefont\small},
  book arrow/.style={figure flow}
}
% 左下 x/y、宽、高、标题、标题带颜色。仅用于已有的阶段或模块分区。
% 使用 x=y=1mm 的坐标；在内容之前绘制，以免遮挡节点与连线。
\newcommand{\figurePanel}[6]{%
  \path[figure panel] (#1,#2) rectangle ({#1+#3},{#2+#4});%
  \begin{scope}
    \clip[rounded corners=2.5mm] (#1,#2) rectangle ({#1+#3},{#2+#4});%
    \fill[#6] (#1,{#2+#4-8}) rectangle ({#1+#3},{#2+#4});%
  \end{scope}
  \draw[draw=FigureStroke,line width=.6pt]
    (#1,{#2+#4-8})--({#1+#3},{#2+#4-8});%
  \node[figure heading] at ({#1+#3/2},{#2+#4-4}) {#5};%
}

% 其他卷保留既有 book node/book arrow；第一卷局部应用科研样式。
\tikzset{book node/.style={draw=none,fill=BookPaper,minimum height=10mm,
  inner sep=6pt,font=\heiti\small,align=center},
  book arrow/.style={-{Stealth[length=2mm]},line width=0.7pt,draw=BookMuted}}

% 神经网络章节的可复用图形语法；颜色在 layout.tex 集中定义。
\tikzset{
  nn neuron/.style={circle,minimum size=8mm,inner sep=0pt,
    line width=1.1pt,text=NNInk,draw=NNHidden,fill=NNHidden!18},
  nn input/.style={nn neuron,draw=NNInput,fill=NNInput!18},
  nn hidden/.style={nn neuron,draw=NNHidden,fill=NNHidden!18},
  nn output/.style={nn neuron,draw=NNOutput,fill=NNOutput!18},
  nn link/.style={draw=NNLine,line width=.85pt},
  nn flow/.style={-{Stealth[length=2mm,width=1.4mm]},
    draw=NNInput,line width=1.2pt},
  nn gradient/.style={-{Stealth[length=2mm,width=1.4mm]},
    draw=NNGradient,line width=1.2pt,dashed},
  nn heading/.style={font=\figurefont\bfseries\footnotesize,text=NNInk},
  nn annotation/.style={font=\figurefont\footnotesize,text=BookMuted},
  nn transform/.style={rounded corners=1.5mm,draw=NNOutput,
    fill=NNOutput!18,line width=1.1pt,minimum height=9mm,
    inner sep=3pt,align=center,text=NNInk}
}

% 操作示例的数值网格：左上 x/y、行数、列数、格宽、边框色、行/列/值。
% 空间索引从零开始，色彩和字号沿用神经网络图形语法。
\newcommand{\cnnGrid}[7]{%
  \draw[draw=#6,line width=.95pt] (#1,#2)
    rectangle ({#1+#4*#5},{#2-#3*#5});%
  \ifnum#4>1\relax
  \pgfmathtruncatemacro{\cnnGridLast}{#4-1}%
  \foreach \cnnGridIndex in {1,...,\cnnGridLast}{%
    \draw[nn link,line width=.5pt] ({#1+\cnnGridIndex*#5},#2)--
      ({#1+\cnnGridIndex*#5},{#2-#3*#5});%
  }%
  \fi
  \ifnum#3>1\relax
  \pgfmathtruncatemacro{\cnnGridLast}{#3-1}%
  \foreach \cnnGridIndex in {1,...,\cnnGridLast}{%
    \draw[nn link,line width=.5pt] (#1,{#2-\cnnGridIndex*#5})--
      ({#1+#4*#5},{#2-\cnnGridIndex*#5});%
  }%
  \fi
  \foreach \cnnGridRow/\cnnGridCol/\cnnGridValue in {#7}{%
    \node at ({#1+(\cnnGridCol+.5)*#5},{#2-(\cnnGridRow+.5)*#5})
      {$\cnnGridValue$};%
  }%
}

% CNN 的长方体表示空间与通道；局部块、角点连线及平面 MLP 共用颜色。
\pgfdeclarelayer{cnn background}
\tikzset{
  cnn plane/.style={draw=NNHidden,line width=1.1pt},
  cnn flow/.style={nn flow},
  cnn skip/.style={nn flow,draw=NNHidden,line width=1.1pt,rounded corners=2mm},
  cnn mapping/.style={draw=NNHidden!48,line width=.85pt,
    dash pattern=on 2.2pt off 1.5pt},
  cnn scale/.style={draw=NNLine,line width=.85pt},
  cnn title/.style={font=\figurefont\bfseries\footnotesize,text=NNInk,align=center},
  cnn size/.style={font=\figurefont\footnotesize,text=NNInk,align=center},
  cnn operation/.style={font=\figurefont\footnotesize,text=BookMuted,align=center},
  cnn neuron/.style={nn hidden,minimum size=3.5mm},
  cnn merge/.style={circle,draw=NNInk,fill=NNCanvas,
    minimum size=5.5mm,inner sep=0pt,line width=1.1pt}
}
% 单行架构的统一标题线、尺寸线；真正并行的分支可保留局部标题。
\newcommand{\cnnArchitecture}[2]{%
  \def\cnnTitleY{#1}\def\cnnSizeY{#2}%
  \pgfsetlayers{cnn background,main}%
}
\newcommand{\cnnLayerTitle}[3]{%
  \ifdefined\cnnTitleY
    \node[cnn title,anchor=north] at ({#1},\cnnTitleY) {#3};%
  \else
    \node[cnn title,anchor=south] at ({#1},{#2}) {#3};%
  \fi
}
\newcommand{\cnnLayerSize}[3]{%
  \ifdefined\cnnSizeY
    \node[cnn size,anchor=north] at ({#1},\cnnSizeY) {#3};%
  \else
    \node[cnn size,anchor=north] at ({#1},{#2}) {#3};%
  \fi
}
% 名称、外框左下 x/y、正面宽/高、示意厚度、颜色、标题、实际形状。
\newcommand{\cnnFeature}[9]{%
  \pgfmathsetmacro{\cnnDepth}{.7*#6}%
  \colorlet{cnn-#1-color}{#7}%
  \draw[cnn plane,draw=#7,fill=#7!12]
    (#2,{#3+#5+\cnnDepth})--({#2+#4},{#3+#5+\cnnDepth})--
    ({#2+#4+\cnnDepth},{#3+#5})--({#2+\cnnDepth},{#3+#5})--cycle;%
  \draw[cnn plane,draw=#7,fill=#7!30]
    (#2,{#3+\cnnDepth})--({#2+\cnnDepth},#3)--
    ({#2+\cnnDepth},{#3+#5})--(#2,{#3+#5+\cnnDepth})--cycle;%
  \draw[cnn plane,draw=#7,fill=#7!20]
    ({#2+\cnnDepth},#3) rectangle ({#2+#4+\cnnDepth},{#3+#5});%
  \coordinate (#1-west) at (#2,{#3+#5/2+\cnnDepth/2});%
  \coordinate (#1-east) at ({#2+#4+\cnnDepth},{#3+#5/2});%
  \coordinate (#1-north) at ({#2+(#4+\cnnDepth)/2},{#3+#5+\cnnDepth});%
  \coordinate (#1-south) at ({#2+(#4+\cnnDepth)/2},#3);%
  % 相邻张量连接右侧面与左侧面；分别保留前后、上下四个角点。
  \foreach \cnnLevel/\cnnY in {bottom/0,top/1}{%
    \coordinate (#1-left-front-\cnnLevel) at ({#2+\cnnDepth},{#3+\cnnY*#5});%
    \coordinate (#1-left-back-\cnnLevel) at (#2,{#3+\cnnY*#5+\cnnDepth});%
    \coordinate (#1-right-front-\cnnLevel) at ({#2+#4+\cnnDepth},{#3+\cnnY*#5});%
    \coordinate (#1-right-back-\cnnLevel) at ({#2+#4},{#3+\cnnY*#5+\cnnDepth});%
  }%
  \foreach \cnnCorner/\cnnX/\cnnY in {sw/0/0,se/1/0,nw/0/1,ne/1/1}{%
    \coordinate (#1-face-\cnnCorner) at ({#2+\cnnDepth+\cnnX*#4},{#3+\cnnY*#5});%
    \coordinate (#1-patch-\cnnCorner) at ({#2+\cnnDepth+(.18+.37*\cnnX)*#4},{#3+(.22+.31*\cnnY)*#5});%
    \coordinate (#1-patch-back-\cnnCorner) at ({#2+(.18+.37*\cnnX)*#4},{#3+(.22+.31*\cnnY)*#5+\cnnDepth});%
    \coordinate (#1-cell-\cnnCorner) at ({#2+\cnnDepth+(.2+.14*\cnnX)*#4},{#3+(.32+.14*\cnnY)*#5});%
    \coordinate (#1-cell-back-\cnnCorner) at ({#2+(.2+.14*\cnnX)*#4},{#3+(.32+.14*\cnnY)*#5+\cnnDepth});%
  }%
  \foreach \cnnUnit/\cnnRatio in {1/.2,2/.5,3/.8}{%
    \coordinate (#1-vout-\cnnUnit) at ({#2+#4+\cnnDepth},{#3+\cnnRatio*#5});%
    \coordinate (#1-vin-\cnnUnit) at ({#2+\cnnDepth},{#3+\cnnRatio*#5});%
  }%
  \cnnLayerTitle{#2+(#4+\cnnDepth)/2}{#3+#5+\cnnDepth+2.5}{#8}%
  \cnnLayerSize{#2+(#4+\cnnDepth)/2}{#3-2.5}{#9}%
}
\newcommand{\cnnScaleLinks}[2]{%
  \begin{pgfonlayer}{cnn background}%
  \foreach \cnnSide in {front,back}{\foreach \cnnLevel in {bottom,top}{%
    \draw[cnn scale] (#1-right-\cnnSide-\cnnLevel)--(#2-left-\cnnSide-\cnnLevel);%
  }%
  }%
  \end{pgfonlayer}%
}
% 小块沿通道厚度贯穿，以所属张量的深色突出。
\newcommand{\cnnLocalBlock}[2][NNHidden]{%
  \draw[draw=#1!80!NNInk,fill=#1!38,line width=1.1pt]
    (#2-back-nw)--(#2-back-ne)--(#2-ne)--(#2-nw)--cycle;%
  \draw[draw=#1!80!NNInk,fill=#1!65,line width=1.1pt]
    (#2-back-sw)--(#2-sw)--(#2-nw)--(#2-back-nw)--cycle;%
  \draw[draw=#1!80!NNInk,fill=#1!48,line width=1.1pt]
    (#2-sw)--(#2-se)--(#2-ne)--(#2-nw)--cycle;%
}
\newcommand{\cnnMap}[2]{%
  % 局部块同样只连接彼此相向的侧面：右侧面的四角对应左侧面的四角。
  \foreach \cnnFrom/\cnnTo in {se/sw,ne/nw,back-se/back-sw,back-ne/back-nw}{%
    \draw[cnn mapping] (#1-patch-\cnnFrom)--(#2-cell-\cnnTo);%
  }%
  \cnnLocalBlock[cnn-#1-color]{#1-patch}%
  \cnnLocalBlock[cnn-#2-color]{#2-cell}%
}
\newcommand{\cnnWholeLinks}[2]{%
  \foreach \cnnFrom in {1,2,3}{\foreach \cnnTo in {1,2,3}{%
    \draw[nn link] (#1-vout-\cnnFrom)--(#2-vin-\cnnTo);%
  }}%
}
% 平面全连接层：可选颜色、名称、中心 x/y、层名、单元数。
\newcommand{\cnnFC}[6][NNHidden]{%
  \draw[draw=#1,fill=#1!8,line width=1.1pt]
    ({#3-3.4},{#4-13}) rectangle ({#3+3.4},{#4+13});%
  \foreach \cnnUnit/\cnnY in {1/-8,2/0,3/8}{%
    \node[cnn neuron,draw=#1,fill=#1!18] (#2-\cnnUnit) at (#3,{#4+\cnnY}) {}; %
    \coordinate (#2-vin-\cnnUnit) at ({#3-1.75},{#4+\cnnY});%
    \coordinate (#2-vout-\cnnUnit) at ({#3+1.75},{#4+\cnnY});%
  }%
  \coordinate (#2-west) at ({#3-3.4},#4);%
  \coordinate (#2-east) at ({#3+3.4},#4);%
  \cnnLayerTitle{#3}{#4+15.5}{#5}%
  \cnnLayerSize{#3}{#4-15.5}{#6}%
}
\newcommand{\cnnDenseLinks}[2]{\cnnWholeLinks{#1}{#2}}

% 平面计算图采用浅色圆角模块、简短算子标签和连贯的圆角跨层路径。
\tikzset{
  nn operator/.style={draw=NNHidden,fill=NNHidden!18,rounded corners=1.5mm,
    line width=1.1pt,minimum height=10mm,inner sep=0pt,align=center},
  nn operator path/.style={nn flow,draw=NNHidden,rounded corners=2.5mm},
  nn pointwise/.style={circle,draw=NNHidden,fill=NNCanvas,line width=1.1pt,
    minimum size=6mm,inner sep=0pt}
}
% 名称、中心 x/y、核尺寸。
\newcommand{\nnConvOp}[4]{%
  \node[nn operator,fill=NNHidden!18,minimum width=10mm,minimum height=13mm]
    (#1) at (#2,#3) {Conv\\$#4$};%
}
\newcommand{\nnReluOp}[3]{%
  \node[nn operator,fill=NNHidden!18,minimum width=10mm] (#1) at (#2,#3) {ReLU};%
}
\newcommand{\nnSigmoidOp}[3]{%
  \node[nn operator,fill=NNHidden!18,minimum width=17mm] (#1) at (#2,#3) {Sigmoid};%
}
\newcommand{\nnBNOp}[3]{%
  \node[nn operator,fill=NNHidden!18,minimum width=8mm] (#1) at (#2,#3) {BN};%
}
% 可选圆直径、名称、中心 x/y；青色圆形 + 表示通道拼接，图注说明其语义。
\newcommand{\nnConcatOp}[4][6mm]{%
  \node[nn pointwise,draw=NNOutput,minimum size=#1]
    (#2) at (#3,#4) {$+$};%
}
\newcommand{\nnTensorState}[5]{%
  \node[draw=#4,fill=#4!8,line width=1.1pt,minimum width=8mm,
    minimum height=20mm,inner sep=0pt] (#1) at (#2,#3) {};%
  \foreach \nnOffset in {-6,0,6}{%
    \node[circle,draw=#4,fill=#4!20,line width=1.1pt,
      minimum size=3mm,inner sep=0pt] at (#2,{#3+\nnOffset}) {};%
  }%
  \node[nn heading,anchor=south] at (#2,{#3+12}) {#5};%
}


% 用于第18章旧图形组件的兼容样式；语义颜色与后续章节相同。
\tikzset{
  nn legacy node/.style={nn operator,inner sep=5pt},
  nn legacy arrow/.style={nn flow,draw=NNWire},
  nn frame/.style={draw=NNLine,fill=NNPanel,line width=.85pt,rounded corners=1.5mm}
}

